% ============================================================================
% Chapter 4 — Performance Poetry Forms
% ============================================================================
\chapter{Performance Poetry Forms}
\label{ch:performance}

Performance poetry (\ar{الشعر الأدائي}) designates forms in which the compositional and performative acts are inseparable — where the context of delivery (a dueling stage, a tribal gathering, a dance circle) is built into the rules of the form itself. This chapter covers four major traditions: improvised poetic dueling (\ar{المحاورة / القلطة}), the Southern Ardah (\ar{العرضة الجنوبية}), the Samri nocturnal tradition (\ar{السامري}), and the Levantine Zajal (\ar{الزجل}).

% ============================================================================
\section{\ar{شعر المحاورة والقلطة} — Improvised Poetic Dueling}
\label{sec:muhawara}

\subsection{Definition and Regional Names}

\ar{المحاورة} (``the conversation'') and \ar{القلطة} (``the advance/challenge'') refer to the \textit{same artistic tradition} under different regional names:

\begin{table}[H]
\centering
\caption{Regional names for improvised poetic dueling}
\label{tab:muhawara_names}
\renewcommand{\arraystretch}{1.4}
\begin{tabularx}{\textwidth}{llX}
\toprule
\textbf{Region} & \textbf{Name Used} & \textbf{Notes} \\
\midrule
Hijaz (Mecca, Medina, Jeddah) & \ar{المحاورة / الرد / الرديح} & \ar{الرديح} is considered the oldest form name; \ar{الرد} is common \\
Najd / Gulf & \ar{القلطة} & From \ar{قلط على} = ``to advance toward someone''; the challenge posture \\
Eastern Arabia & \ar{الرديّة} & Also used: \ar{القلطة} \\
Tihama coast & \ar{الرد} & Generic; also improvised \\
\bottomrule
\end{tabularx}
\end{table}

\subsection{The Core Mechanics}

\ar{المحاورة / القلطة} is improvised, metered, rhymed poetry recited as a \textit{real-time duel} between two poets before a live audience. The defining characteristics:

\begin{enumerate}
  \item \textbf{Complete improvisation}: No pre-writing or preparation is permitted. The challenge to the tradition's integrity is the \textit{spontaneity} of composition. A poet who has written their verses in advance has violated the spirit of the form.
  \item \textbf{Musical delivery}: Verses are not recited but \textit{sung} to a specific melodic path (\ar{طريق}).
  \item \textbf{Audience participation}: Audience members (or a chorus group, \ar{الصفوف}) repeat verses or provide rhythmic support.
\end{enumerate}

\subsection{The Tariq System (\ar{نظام الطريق})}

In \ar{المحاورة}, the meter is not specified as a Khalilian or Nabati bahr but as a \ar{طريق} (``path'' or ``melody-meter''). The tariq determines:
\begin{itemize}
  \item The prosodic foot structure
  \item The melody used for singing
  \item The pace and rhythm of the exchange
\end{itemize}

Both poets are \textit{bound to the same tariq} for the duration of the exchange. A poet who breaks the tariq loses the round.

\subsection{The Opening and Response Structure}

\begin{description}
  \item[\ar{الفتح / الفتل}] \textbf{The opening move}: The first poet introduces an argument, image, or challenge in verse. This must be metrically and melodically correct.
  \item[\ar{النقض / الرد}] \textbf{The response}: The second poet must \textit{immediately} counter — accepting or refuting the argument in metrically matching verse. The quality of response is judged on: speed, prosodic accuracy, the cleverness of the counter-argument, and the power of the imagery.
\end{description}

\subsection{Evaluation Criteria}

A \ar{محاورة} exchange is judged — by critics, by fellow poets, and by the audience — on:
\begin{enumerate}
  \item \textbf{Prosodic accuracy}: Both poets must maintain meter and rhyme precisely.
  \item \textbf{Speed of response}: Hesitation is a sign of weakness.
  \item \textbf{Argument quality}: The response must \textit{engage} with the specific claim or image of the preceding verse — not just add a new independent verse.
  \item \textbf{Imagery}: Striking metaphors and unexpected comparisons score highly.
  \item \textbf{Audience reaction}: The audience's spontaneous response (applause, repetition of the verse) is itself a form of judgment.
\end{enumerate}

\subsection{Types of \ar{محاورة}}

\begin{description}
  \item[\ar{الفاضلة}] \textit{Noble/elevated exchange}: Addresses high literary, ethical, or social themes. More dignified tone.
  \item[\ar{الفضولية}] \textit{Tribal/popular exchange}: Involves direct personal and tribal challenge; more aggressive.
\end{description}

\begin{rulebox}{AI Generation Rules for \ar{المحاورة}}
\begin{enumerate}
  \item The system must simulate the exchange: given the first poet's verse, generate a valid counter-verse.
  \item The counter-verse \textbf{must use the same tariq (meter + rhyme)} as the opening verse.
  \item The counter-verse must \textbf{engage with the specific argument} of the opening — not simply add a new theme.
  \item Language must be in the appropriate regional dialect (Hijazi for \ar{محاورة}, Najdi for \ar{قلطة}).
  \item The response should exhibit \ar{براعة الارتجال}: wit, speed of thought, and surprise.
  \item For a full simulated exchange: generate multiple rounds (minimum 3 exchanges) maintaining prosodic consistency throughout.
\end{enumerate}
\end{rulebox}

% ============================================================================
\section{\ar{العرضة الجنوبية} — The Southern Ardah}
\label{sec:ardah}

\subsection{Definition and Origins}

Al-Ardah al-Janubiyya (\ar{العرضة الجنوبية}) is a martial performance poetry tradition from the southern Saudi regions: Asir (\ar{عسير}), Jizan (\ar{جازان}), Najran (\ar{نجران}), and Al-Bahah (\ar{الباحة}). At least four centuries old, it originated as a martial display (\ar{عرض عسكري}) designed to mobilise fighters and demonstrate tribal unity before battle.

\subsection{The Performance Structure}

The Ardah is a precisely choreographed performance:
\begin{enumerate}
  \item \textbf{Formation}: Participants stand in \textbf{two parallel rows facing each other}, with approximately two metres between the rows.
  \item \textbf{The central poet}: One or more poets (\ar{شاعر العرضة}) stand in the \textit{space between the two rows}.
  \item \textbf{Delivery}: The poet turns to the first row and recites the \textit{sadr} (first hemistich) of a verse; then turns to the second row and delivers the \textit{ajuz} (second hemistich).
  \item \textbf{Collective chant}: Both rows repeat the verse in unison (\ar{الإنشاد الجماعي}).
  \item \textbf{Multi-poet exchange}: In elaborate performances, multiple poets (sometimes four) engage in \ar{محاورة} between the rows.
  \item \textbf{Instruments}: The \ar{زِير} (a percussion instrument) provides the rhythmic pulse; sometimes drums (\ar{طبول}) are added.
  \item \textbf{Props}: Participants traditionally carry swords, daggers, or rifles.
\end{enumerate}

\subsection{Poetic Characteristics of the Ardah}

\begin{itemize}
  \item Poems are typically \textbf{short} — often not exceeding ten verses.
  \item Language is in the \textbf{colloquial Arabic of the southern regions} (Asiri, Jizan dialects).
  \item The primary meter historically was rhymed stanzas (\ar{أبيات مقفاة}). Over two centuries ago, poets transitioned to what is known as \ar{شعر الشقر} (a specific southern regional form with its own rules).
  \item \textbf{Themes}: tribal pride (\ar{فخر}), martial enthusiasm (\ar{حماسة}), heroism (\ar{بطولة}), and historical tribal epics.
\end{itemize}

\subsection{Historical Forms of the Ardah}

Three historical types of Ardah existed:
\begin{description}
  \item[\ar{عرضة الحرب}] \textit{War Ardah}: Performed before battle; specific drum rhythm; only warriors going into battle participate.
  \item[\ar{عرضة الخروج}] \textit{Departure Ardah}: Performed in a straight line, two rows moving forward; short poetic phrases.
  \item[\ar{عرضة الدخول}] \textit{Entrance Ardah}: Performed entering a village; used for welcoming guests or celebrating returning warriors.
\end{description}

\begin{rulebox}{AI Generation Rules for the Ardah}
\begin{enumerate}
  \item The system must know the \textbf{occasion} (battle eve, tribal celebration, historical performance) to set the correct tone.
  \item Language must be in the \textbf{southern Saudi dialects} (Asiri/Jizan/Najrani), not Najdi or Gulf Arabic.
  \item Poems should be short (4–10 verses) and each verse must be \textit{self-contained in meaning} to enable collective repetition.
  \item The sadr/ajuz structure must be particularly clear — the two rows must be able to grasp and repeat each hemistich immediately.
  \item Themes: martial pride (\ar{حماسة}) and tribal solidarity (\ar{وحدة قبلية}) are primary; love poetry is inappropriate for Ardah.
  \item If generating \ar{شعر الشقر}, apply the specific meter and rhyme patterns of that form.
\end{enumerate}
\end{rulebox}

% ============================================================================
\section{\ar{السامري} — The Samri Tradition}
\label{sec:samri}

\subsection{Definition and Origins}

Samri (\ar{السامري}) is a folkloric music, poetry, and dance tradition from Najd, Saudi Arabia, approximately 300 years old. Its name derives from \ar{السمر} (staying up late at night in conversation and entertainment) — it is intrinsically a \textbf{nocturnal} art form, performed at evening and night-time gatherings.

\subsection{The Performance Structure}

\begin{enumerate}
  \item \textbf{Formation}: Two opposing rows of men face each other, seated on their knees.
  \item \textbf{Synchronized movement}: Men sway, clap, and sing in synchrony — the bodily movement is as important as the verse.
  \item \textbf{Al-Muallim (المعلم)}: A master teacher-vocalist leads the group by \textit{teaching successive verses in real time}. This is the Samri equivalent of the \ar{شاعر العرضة}.
  \item \textbf{Call-and-response}: One group \textit{opens} (\ar{يشيل}) with a verse; the opposing group gives the \textit{mard} (\ar{المرد}) — repetition/response.
  \item \textbf{Opening chants}: The performance begins with \ar{الوانين} (opening chants) establishing the melodic foundation.
\end{enumerate}

\subsection{Poetic Content}

\begin{itemize}
  \item Samri is built on \textbf{Nabati (colloquial) poetry}, specifically using the \textbf{Mashhub meter} with its dual-rhyme system.
  \item \ar{بحر السامري} is also recognised as a distinct meter in the Nabati system.
  \item The poetry centres on works by foundational poets: Muhammad ibn Labun and Muhsin al-Hazani.
  \item Muhsin al-Hazani is credited with introducing the Mashhub meter to Samri poetry and formalising the dual-rhyme system for musical performance.
  \item Contemporary Samri performances include modern Nabati poetry on personal and emotional themes.
\end{itemize}

\subsection{Instruments}

\begin{table}[H]
\centering
\caption{Samri instruments}
\renewcommand{\arraystretch}{1.4}
\begin{tabularx}{\textwidth}{llX}
\toprule
\textbf{Instrument} & \textbf{Arabic} & \textbf{Role} \\
\midrule
Al-Tar / Al-Tabl & \ar{الطار / الطبل} & Main hand drum; provides the rhythmic pulse \\
Al-Mirjaf & \ar{المرجاف} & Large frame drum; produces deep tones \\
Al-Mard & \ar{المرد} & Small response drum \\
Al-Daf & \ar{الدف} & Standard frame drum \\
\bottomrule
\end{tabularx}
\end{table}

\subsection{Regional Samri Styles}

\begin{table}[H]
\centering
\caption{Regional Samri variations}
\renewcommand{\arraystretch}{1.4}
\begin{tabularx}{\textwidth}{lX}
\toprule
\textbf{Region} & \textbf{Characteristics} \\
\midrule
Aniza / Qassim & Heaviest, most technically demanding; considered the canonical form \\
Duwasir (southern Najd) & Both heavy and light variants; faster tempos \\
Coastal regions & Heavier percussion, slower pacing \\
Hail & Emphasises poetic content over musical elaboration \\
\bottomrule
\end{tabularx}
\end{table}

\begin{rulebox}{AI Generation Rules for Samri Poetry}
\begin{enumerate}
  \item Samri poetry uses the \textbf{Mashhub meter} — the dual-rhyme constraint is mandatory.
  \item Verses must be \textbf{short and rhythmically clear} for group chanting — long complex sentences are inappropriate.
  \item Language: Central/Northern Najdi dialect.
  \item Themes: Love, longing, tribal camaraderie, nocturnal gathering mood — not war or formal praise.
  \item The \ar{وانين} opening is a distinct structural element: a melodic phrase without full lexical content that sets the musical key.
  \item Each verse must be \textit{complete and repeatable} — the opposite row must be able to chant it without knowing the next verse.
\end{enumerate}
\end{rulebox}

% ============================================================================
\section{\ar{الزجل} — The Zajal}
\label{sec:zajal}

\subsection{Definition and Origins}

Zajal (\ar{الزجل}) is colloquial, strophic, semi-improvised oral poetry originating in al-Andalus and still vigorously alive in Lebanon, Syria, Palestine, and Egypt. The earliest documented Zajal poet was Ibn Quzman of Cordoba (1078–1160). The form reached the Levant and Egypt in the 13th–14th centuries following the Reconquista.

The name derives from \ar{زَجَلَ} — to raise the voice in song or to shout joyfully.

\subsection{Structural Characteristics}

Unlike classical poetry's monorhyme structure, Zajal is \textbf{strophic}: each stanza has an internal rhyme pattern. The most common Zajal structure is the \textbf{qasida zajal} form:

\begin{itemize}
  \item \textbf{Al-Matla'} (\ar{المطلع}): Opening couplet establishing the main rhyme (AA).
  \item \textbf{Al-Qafla} (\ar{القفلة}): A refrain (the same rhyme as the matla') that recurs at the end of each stanza.
  \item \textbf{Al-Dour} (\ar{الدور}): The body stanzas, each with an internal rhyme (BB, CC...) before returning to the qafla.
\end{itemize}

A stanza of three dour lines + one qafla line = the \textbf{murabba'} form (4-line strophe), the most common Zajal unit.

\subsection{Lebanese Zajal Forms}

A Lebanese Zajal evening (\ar{ليلة زجلية}) typically progresses through:
\begin{description}
  \item[\ar{القصيدة}] Formal ode recitation to open the evening.
  \item[\ar{المعنّى والقريدي}] Popular debate structures; the dominant form. Includes sub-variants like \ar{الموخمّس المردود}.
  \item[\ar{الغزل}] Love-lyric forms including \ar{الموشحات} variants — joyful, flirtatious.
  \item[\ar{الشروقي}] Concluding love-lament form.
\end{description}

\subsection{The Zajal Dueling Format}

Lebanese Zajal is primarily a \textbf{competitive duel} (\ar{مجادلة زجلية}) between two or more \ar{زجّالين} (Zajal poets):
\begin{enumerate}
  \item Poets improvise verses on the same topic and in the same form.
  \item A crucial rule: when a poet responds, they \textbf{must begin with the last word} of the previous poet's stanza.
  \item The exchange is accompanied by percussion (tambourines, \ar{رقّ}) and sometimes a ney flute.
  \item A chanting chorus (\ar{الكورس}) repeats verses and refrains.
\end{enumerate}

\subsection{Metrical System}

Zajal uses two distinct metrical systems in parallel:
\begin{description}
  \item[\textbf{Quantitative}] Based on Khalilian classical meters (e.g., \ar{المعنّى} and some forms scan according to Sari', Rajaz, Wafir).
  \item[\textbf{Stress-syllabic}] Based on Syriac/Aramaic metrics (e.g., many sub-forms of the \ar{القريدي}).
\end{description}

\begin{rulebox}{AI Generation Rules for Zajal}
\begin{enumerate}
  \item The system must identify the \textbf{specific Zajal form} (Lebanese, Egyptian, Moroccan) and the \textbf{dialect} to be used.
  \item Apply the \textbf{strophic structure}: matla' → dour stanzas → recurring qafla.
  \item For dueling format: the response poet's stanza \textit{must begin with the last word of the previous stanza}.
  \item Language must be the appropriate colloquial dialect (Lebanese Arabic for Lebanese Zajal, Egyptian for Egyptian, etc.).
  \item The matla' rhyme is the binding constant of the entire poem — it must be maintained precisely.
  \item Zajal permits stress-syllabic meters not found in Khalilian prosody — do not attempt to validate Zajal with the classical bohour library.
\end{enumerate}
\end{rulebox}

% ============================================================================
\section{Other Performance Forms}
\label{sec:other_performance}

\subsection{\ar{الحداء} — The Camel-Driver's Call}

Al-Hida' (\ar{الحداء}) is one of the oldest Arabic poetic forms — verse composed and sung during camel travel to set the pace and sustain morale. It uses the Rajaz meter (or its Nabati equivalent, Bahr al-Hida'). Characteristics:
\begin{itemize}
  \item Short, simple, highly rhythmic — designed to synchronise with the camel's gait.
  \item Often improvised on the spot.
  \item Themes: encouragement, journey, longing for the destination.
\end{itemize}

\subsection{\ar{الموال / المواليا} — The Mawwal}

The mawwal (\ar{الموال}) is a popular solo singing form in Levantine and Egyptian traditions, typically in 4 or 7 lines with a complex rhyme scheme that incorporates wordplay (\ar{جناس}). The Nabati equivalent was introduced c. 1411 AH with a 3-line stanza featuring melodic complexity.

% ============================================================================
\clearpage
\langsep{}

% ============================================================================
%  ARABIC VERSION
% ============================================================================
\begin{otherlanguage}{arabic}

\section*{أشكال الشعر الأدائي}

\subsection*{شعر المحاورة والقلطة}

المحاورة (في الحجاز) والقلطة (في نجد والخليج) اسمان لفن واحد: هو الشعر الارتجالي التنافسي بين شاعرَين أمام جمهور حيّ. يقوم الشاعران بتبادل الأبيات المرتجلة على نفس الطريق (الوزن والنغم) — وهو أشبه بمبارزة فكرية شعرية.

\textbf{قواعد المحاورة:}
\begin{enumerate}
  \item الارتجال التام: يُحرَّم الاستعداد المسبق؛ القيمة الفنية تكمن في العفوية.
  \item الالتزام بالطريق (الوزن والنغم) طوال المبارزة.
  \item كل بيت إجابة يجب أن يُعالج حجة الشاعر السابق، لا مجرد إضافة موضوع جديد.
  \item من لم يستطع الرد سريعًا هُزم.
\end{enumerate}

\textbf{قواعد الذكاء الاصطناعي:} لتوليد رد المحاورة، يجب على النظام: (1) الحفاظ على نفس الطريق، (2) معالجة الحجة الشعرية السابقة مباشرة، (3) استخدام اللهجة المحلية المناسبة (الحجازية أو النجدية).

\subsection*{العرضة الجنوبية}

العرضة الجنوبية فن أدائي من مناطق عسير وجازان ونجران والباحة. يقف المشاركون في صفَّين متقابلَين، ويقف الشاعر في المنتصف بين الصفين يُلقي صدر البيت تجاه الصف الأول وعجزه تجاه الصف الثاني، ثم يُرددان معًا.

\textbf{خصائص شعر العرضة:}
\begin{itemize}
  \item قصائد قصيرة (4–10 أبيات عادةً).
  \item اللغة من لهجات المنطقة الجنوبية (العسيرية والجيزانية).
  \item الموضوعات: الفخر والحماسة والبطولة والتاريخ القبلي.
  \item كل بيت يجب أن يكون مكتفيًا بنفسه لأن الجمهور يُردده.
\end{itemize}

\subsection*{السامري}

السامري فن ليلي نجدي عمره نحو ثلاثمئة عام. يجلس الرجال في صفَّين متقابلَين، يتمايلون ويصفقون وينشدون الشعر النبطي (في الغالب على بحر المسحوب) مع المعلم الذي يُعلّمهم الأبيات أولًا بأول. أبرز شعراء السامري: محمد بن لعبون ومحسن الحزاني الذي وضع نظام القافية المزدوجة في السامري.

\textbf{قواعد الذكاء الاصطناعي للسامري:}
\begin{itemize}
  \item استخدام بحر المسحوب مع القافية المزدوجة.
  \item أبيات قصيرة واضحة الإيقاع لتمكين الإنشاد الجماعي.
  \item اللهجة النجدية الوسطى/الشمالية.
  \item الموضوعات: الحب والحنين والسمر والرفقة — لا الحرب ولا المديح الرسمي.
\end{itemize}

\subsection*{الزجل}

الزجل فن الشعر العامي المتقفى في البلاد الشامية ومصر والمغرب. أقدم زجّال موثق هو ابن قزمان الأندلسي (1078–1160م). في الزجل اللبناني، يتنافس شاعران أو أكثر بارتجال الأبيات على نفس الموضوع. القاعدة الذهبية: الشاعر المجيب يجب أن يبدأ مقطعه بآخر كلمة في مقطع الشاعر السابق.

\textbf{الشكل البنائي للزجل:}
\begin{itemize}
  \item \textbf{المطلع}: بيتان يُثبتان القافية الأساسية (أأ).
  \item \textbf{القفلة}: لازمة تعود في نهاية كل مقطع على نفس قافية المطلع.
  \item \textbf{الدور}: مقاطع الجسم، لكل مقطع قافية داخلية خاصة (بب، جج...) قبل العودة للقفلة.
\end{itemize}

\end{otherlanguage}
